% PR #1113 的 MWE：#1108 与 #1110。
% 用 xelatex 编译。修复前：第一节报 Improper \spacefactor，第二节每组上下两行
% 不等宽（链接写法比直接输入多一枚 \CJKecglue）；修复后不报错，每组两行等宽。
\documentclass{article}
\usepackage{xeCJK}
\usepackage{hyperref}
\setCJKmainfont{FandolSong-Regular.otf}
% 正文里的 \par 要放在宏里，直接写进 \href 的参数会报参数扫描错误（与本问题无关）。
\newcommand{\EndPar}{Hello\par}
\newcommand{\NewPar}{Hello\par\leavevmode}
\newcommand{\TwoPar}{a\par b\par\leavevmode}
\newcommand{\BoxPar}{a\par\vbox{\hsize=3cm b\par}\leavevmode c\par\leavevmode}
\begin{document}
\section*{\#1108：链接正文结束段落}
\href{https://example.com}{\EndPar}
后

\section*{\#1110：链接正文结束段落后又开始新段落}
每组第一行是链接写法，第二行是删去 \verb|\href| 的直接输入，最后一行应等宽。

\medskip
\href{https://example.com}{\NewPar}后\par
Hello\par\leavevmode 后

\medskip
\href{https://example.com}{\TwoPar}后\par
a\par b\par\leavevmode 后

\medskip
\href{https://example.com}{\BoxPar}后\par
a\par\vbox{\hsize=3cm b\par}\leavevmode c\par\leavevmode 后
\end{document}
